\documentclass[11pt]{article}

% Essential packages
\usepackage[margin=1in]{geometry}
\usepackage{amsmath,amssymb,amsthm}
\usepackage{mathtools}
\usepackage[numbers,square]{natbib}
\usepackage{booktabs}
\usepackage{hyperref}
\usepackage{cleveref}
\usepackage{algorithm}
\usepackage{algpseudocode}

% Theorem environments
\theoremstyle{definition}
\newtheorem{definition}{Definition}[section]

\theoremstyle{plain}
\newtheorem{theorem}{Theorem}[section]
\newtheorem{lemma}[theorem]{Lemma}
\newtheorem{corollary}[theorem]{Corollary}
\newtheorem{proposition}[theorem]{Proposition}

\theoremstyle{remark}
\newtheorem{remark}{Remark}[section]
\newtheorem{example}{Example}[section]

% Notation
\newcommand{\fhat}{\hat{f}}
\newcommand{\ghat}{\hat{g}}
\newcommand{\enc}{\mathrm{enc}}
\newcommand{\dec}{\mathrm{dec}}
\newcommand{\im}{\mathrm{Im}}
\newcommand{\TV}{d_{\mathrm{TV}}}
\DeclareMathOperator{\Uniform}{Uniform}

\title{Cipher Maps: Total Functions as Trapdoor Approximations}
\author{Alexander Towell\\\texttt{lex@metafunctor.com}}
\date{\today}

\begin{document}

\maketitle

\begin{abstract}
A cipher map is a total function on bit strings that approximates a latent
function through a one-way trapdoor.  We define four properties---totality,
representation uniformity, correctness, composability---and show how they
arise from three orthogonal construction layers (undefined injection, noise
closure, multiple representations).  We instantiate the abstraction against
three concrete constructions: the HashSet (a space-optimal set membership
structure), the entropy map (achieving Shannon-entropy-optimal value encoding),
and the trapdoor Boolean algebra (an online construction supporting set
operations via bitwise logic).  We derive the composition theorem with proof,
formalize the trusted/untrusted machine model, and present the Singular Hash
Map---a construction achieving the information-theoretic lower bound of
$-\log_2 \varepsilon + H(Y)$ bits per element.
\end{abstract}


%% ============================================================
\section{Introduction}
\label{sec:intro}

Consider a setting in which a trusted party wishes to outsource computation
to an untrusted machine.  The trusted party holds a function
$f : X \to Y$---a lookup table, a set membership predicate, an index---and
wants the untrusted machine to evaluate $f$ on encoded inputs without
learning $f$, $X$, or $Y$.

The central object of this paper is the \emph{cipher map}: a total function
$\fhat : \{0,1\}^n \to \{0,1\}^n$ that the untrusted machine evaluates
blindly.  The key insight is that $\fhat$ is total---every $n$-bit input
produces an $n$-bit output, with no concept of ``in-domain'' or
``out-of-domain.''  The notions of domain, correctness, false positive, and
false negative exist only on the trusted machine, which holds the encoding
and decoding functions (the trapdoor).  The untrusted machine sees only
opaque bit strings flowing through opaque total functions.

This paper develops the theory of cipher maps in a self-contained way.
We define four properties that characterize a cipher map
(\S\ref{sec:properties}), instantiate them against three concrete
constructions from distinct domains (\S\ref{sec:constructions}), formalize
the trusted/untrusted machine model (\S\ref{sec:trust-model}), present the
Singular Hash Map achieving space-optimal performance
(\S\ref{sec:singular-hash}), and prove the composition theorem governing
error accumulation (\S\ref{sec:composition}).

\paragraph{What this is not.}
Cipher maps are not oblivious RAM~\cite{goldreich1996software}: ORAM
reshuffles storage to hide access patterns with polylogarithmic overhead,
while cipher maps use static total functions with noise injection.  Cipher
maps are not fully homomorphic encryption~\cite{gentry2009fully}: FHE
preserves exact algebraic structure on ciphertexts, while cipher maps are
approximate and hash-based.  The closest relative is garbled
circuits~\cite{yao1982protocols}---both use encrypted lookup tables---but
garbled circuits are one-time use and exact, while cipher maps are reusable
and approximate with noise closure.


%% ============================================================
\section{The Cipher Map Abstraction}
\label{sec:abstraction}

\subsection{Definition}

\begin{definition}[Cipher map]
\label{def:cipher-map}
Let $f : X \to Y$ be a \emph{latent function}.  A \emph{cipher map} for
$f$ is a tuple $(\fhat, \enc, \dec, s)$ where:
\begin{enumerate}
\item $\fhat : \{0,1\}^n \to \{0,1\}^n$ is a \textbf{total function} on
      $n$-bit strings;
% FIXED: K -> K(x) to match usage elsewhere (multiplicity varies by element)
\item $\enc : X \times \{0, \ldots, K(x){-}1\} \to \{0,1\}^n$ is an encoding
      function, where $K(x) \geq 1$ is the number of encodings for element $x$;
\item $\dec : \{0,1\}^n \to Y \cup \{\bot\}$ is a decoding function;
\item $s$ is a secret (the trapdoor) from which $\fhat$, $\enc$, and $\dec$
      are derived.
\end{enumerate}
The untrusted machine holds $\fhat$.  The trusted machine holds $\enc$,
$\dec$, and $s$.
\end{definition}

\begin{remark}[Values as constant functions]
\label{rem:values-as-functions}
A cipher value $\enc(v, k)$ can be viewed as a cipher map for the constant
function $f_v(x) = v$ for all $x$.  Under this view, cipher values and
cipher maps are the same kind of object---a total function on bit
strings---simplifying the untrusted machine's interface.
\end{remark}

\subsection{Two Construction Strategies}

Cipher maps arise from two fundamentally different construction strategies:

\begin{description}
\item[Batch construction.]  The trusted machine invests computation upfront
    to find a seed $s$ such that $\fhat$ satisfies correctness constraints
    for all $x \in X$ simultaneously.  This gives tunable correctness
    parameter $\eta$ and optimal space, but requires knowing $X$ and $f$ at
    construction time.  Examples: HashSet (\S\ref{subsec:hashset}), entropy
    map (\S\ref{subsec:entropy-map}).

\item[Online construction.]  The cipher map $\fhat$ is defined directly by
    the hash structure---no seed search is needed.  The secret $s$ is a hash
    key, not a searched seed.  Operations are limited to those the hash
    structure supports algebraically (e.g., set union, intersection).  Space
    and accuracy trade-offs are fixed by the hash width~$n$.  Example:
    trapdoor Boolean algebra (\S\ref{subsec:trapdoor-bool}).
\end{description}

Both strategies produce objects satisfying Definition~\ref{def:cipher-map}.

\subsection{Construction Layers}
\label{subsec:layers}

The four properties (defined in \S\ref{sec:properties}) arise from three
orthogonal type transformations applied to the latent function.

\paragraph{Layer 1: Undefined injection.}
Extend the domain $X$ to $X \cup \{\bot_1, \ldots, \bot_N\}$ by adding $N$
undefined elements.  A function $f : X \to Y$ lifts to a partial function
$\tilde{f}$ that is undefined on the $\bot_i$.  Real elements become a
fraction $|X| / (|X| + N)$ of the extended type.  This controls the space
parameter $\varepsilon$ (the probability that random bits form a valid
codeword).

\paragraph{Layer 2: Noise closure.}
Lift the partial function $\tilde{f}$ to a total function $\fhat$ by
mapping undefined inputs to random outputs via cryptographic hash.  The
untrusted machine cannot distinguish real inputs from undefined
inputs---both produce $n$-bit output.  This yields \textbf{Property~1
(Totality)}.

\paragraph{Layer 3: Multiple representations.}
Replace each value $x$ with $K(x) \geq 1$ distinct representations, keyed
by the secret $s$.  Representational equality implies value equality, but
value equality does not imply representational equality.  Assigning
$K(x) \propto 1/D(x)$ equalizes frequencies.  This yields
\textbf{Property~2 (Representation Uniformity)}.

\medskip

The combined type is
$\mathrm{cipher}(\mathrm{noise}(\mathrm{undef}(X)))$---an element with
multiple representations, total evaluation, and diluted domain.
Properties~3 (Correctness) and~4 (Composability) are emergent: correctness
depends on the construction method; composability depends on the totality
guarantee enabling chained evaluation.

\begin{remark}[Layers as conceptual scaffolding]
We present these as conceptual layers explaining \emph{why} each property
exists, not as algebraic monads.  Whether the layers satisfy formal monad
laws in the approximate setting is an open question that we do not need to
resolve.
\end{remark}


%% ============================================================
\section{Four Properties}
\label{sec:properties}

A cipher map $(\fhat, \enc, \dec, s)$ for $f : X \to Y$ is characterized by
four properties parameterized by $(\eta, \varepsilon, \mu, \delta)$.

\subsection{Property 1: Totality}

\begin{definition}[Totality]
\label{def:totality}
The function $\fhat : \{0,1\}^n \to \{0,1\}^n$ is \emph{total}: for every
$c \in \{0,1\}^n$, $\fhat(c) \in \{0,1\}^n$.  For $c$ chosen uniformly at
random from $\{0,1\}^n \setminus \im(\enc)$, the distribution of $\fhat(c)$
is uniform over $\{0,1\}^n$ under the random oracle model.
\end{definition}

Totality ensures that the untrusted machine cannot distinguish a real query
$\enc(x, k)$ from a random filler string $c \gets \{0,1\}^n$, because both
produce $n$-bit outputs.  If out-of-domain queries returned an error or
$\bot$, the untrusted machine could identify real queries by observing which
inputs produce valid output.

\subsection{Property 2: Representation Uniformity ($\delta$-bounded)}

\begin{definition}[Representation uniformity]
\label{def:rep-uniformity}
Let $D$ be a distribution on $X$.  Define the induced distribution on cipher
values:
\begin{equation}
Q(c) = \sum_{x \in X} D(x) \cdot \frac{|\{k : \enc(x,k) = c\}|}{K(x)}.
\end{equation}
The cipher map has \emph{$\delta$-representation uniformity} if
\begin{equation}
\TV(Q,\, \Uniform(\{0,1\}^n)) \leq \delta,
\end{equation}
where $\TV$ denotes total variation distance.
\end{definition}

Total variation distance has the operational interpretation that no
statistical test can distinguish $Q$ from uniform with advantage exceeding
$\delta$.  To achieve small $\delta$, assign $K(x) \propto 1/D(x)$
encodings per value (homophonic substitution), so that frequent values get
more representations.

\begin{remark}[Marginal uniformity only]
\label{rem:marginal-only}
Representation uniformity bounds the \emph{marginal} distribution of
individual cipher values.  It says nothing about \emph{joint} distributions:
an adversary observing sequences of cipher values can still detect
correlations.  See \S\ref{subsec:granularity} for the encoding granularity
trade-off.
\end{remark}

\subsection{Property 3: Correctness ($\eta$-bounded)}

\begin{definition}[Correctness]
\label{def:correctness}
For $x$ chosen uniformly from $X$ and $k$ chosen uniformly from
$\{0, \ldots, K(x)-1\}$:
\begin{equation}
\Pr_{x,k}\bigl[\dec(\fhat(\enc(x,k))) \neq f(x)\bigr] \leq \eta.
\end{equation}
\end{definition}

The parameter $\eta \in [0,1)$ is a \emph{construction} parameter: it
reflects how many constraints the seed $s$ must satisfy.  Tolerating
$\eta > 0$ makes the seed search faster and the structure smaller.  Once
built, the failing elements are deterministic for a given seed---the same
elements always fail.

\begin{remark}[Approximation as privacy]
Nonzero $\eta$ also provides privacy: a false positive creates plausible
deniability, since the untrusted machine cannot distinguish ``element is in
the set'' from ``cipher map erred.''  With $\eta = 0$, the answer is
deterministic; with $\eta > 0$, each answer carries ambiguity.
\end{remark}

\subsection{Property 4: Composability}

\begin{definition}[Composability]
\label{def:composability}
Let $\fhat$ be a cipher map for $f : X \to Y$ with correctness $\eta_f$,
and $\ghat$ a cipher map for $g : Y \to Z$ with correctness $\eta_g$.
Define $\ghat \circ \fhat : \{0,1\}^n \to \{0,1\}^n$ by
$(\ghat \circ \fhat)(c) = \ghat(\fhat(c))$.  Then $\ghat \circ \fhat$ is a
cipher map for $g \circ f$ with correctness:
\begin{equation}
\label{eq:composition}
\eta_{g \circ f} = 1 - (1-\eta_f)(1-\eta_g) = \eta_f + \eta_g - \eta_f \eta_g.
\end{equation}
\end{definition}

Composability is enabled by totality: since $\fhat$ is total, its output is
always a valid $n$-bit input to $\ghat$.  The untrusted machine can chain
cipher maps without needing to decode intermediate results.  The proof of
\eqref{eq:composition} is given in \S\ref{sec:composition}.

\begin{remark}[Noise under composition]
Noise input to $\fhat$ produces random output (by totality).  With
probability $\varepsilon$, that random output happens to be a valid codeword
for $\ghat$.  We do not try to control noise output---the cipher map is just
a hash.  This is a feature: valid-looking outputs do not prove valid inputs.
\end{remark}

\subsection{Parameter Decomposition}

\begin{table}[ht]
\centering
\caption{Cipher map parameters.}
\label{tab:params}
\begin{tabular}{@{}llll@{}}
\toprule
Parameter & Range & Controls & Determined by \\
\midrule
$\eta$ (correctness) & $[0,1)$ & Fraction of wrong answers &
    Construction (seed search) \\
$\varepsilon$ (noise decode) & $(0,1)$ & Pr[random bits valid] &
    Encoding scheme \\
$\mu$ (value cost) & $(0,\infty)$ & Bits per element for values &
    $\mu = H(Y)$ \\
$\delta$ (uniformity) & $[0,1]$ & TV distance from uniform &
    Multiplicity $K(x)$ \\
$K(x)$ (multiplicity) & $\{1,2,\ldots\}$ & Encodings per element &
    Set to $\propto 1/D(x)$ \\
% FIXED: Added entanglement parameter p (discussed in Section 8, referenced in open questions)
$p$ (entanglement) & $\{1, \ldots, k\}$ & Encoding granularity &
    Correlation structure (\S\ref{subsec:granularity}) \\
\bottomrule
\end{tabular}
\end{table}

The space per element decomposes as:
\begin{equation}
\label{eq:space}
\text{bits/element} = -\log_2 \varepsilon + \mu = -\log_2 \varepsilon + H(Y),
\end{equation}
where $-\log_2 \varepsilon$ bits distinguish signal from noise and $H(Y)$
bits encode the function value.  This decomposition is
information-theoretic.


%% ============================================================
\section{Concrete Constructions}
\label{sec:constructions}

We instantiate the cipher map abstraction against three constructions.  For
each, we identify which of the four properties are satisfied and with what
parameter values.

\subsection{HashSet}
\label{subsec:hashset}

The HashSet is the simplest cipher map: an approximate set membership
structure derived from a seed search.

\paragraph{Latent function.}
Set membership indicator: $f = \mathbf{1}_A : X \to \{0,1\}$ where
$A \subseteq X$.

\paragraph{Construction.}
Choose a seed $s$ such that
\begin{equation}
h(x \| s) = 0^n \quad \text{for all } x \in A,
\end{equation}
where $h : \{0,1\}^* \to \{0,1\}^n$ is a cryptographic hash modeled as a
random oracle.  Membership test: declare $x \in A$ if and only if
$h(x \| s) = 0^n$.

\paragraph{Mapping to cipher map abstraction.}
\begin{itemize}
\item $\fhat(c) = h(c \| s)$, a total function on all bit strings.
\item $\enc(x, 0) = x$ (single encoding per element; $K(x)=1$).
\item $\dec(r) = 1$ if $r = 0^n$, else $0$.
\end{itemize}

\paragraph{Parameter instantiation.}
\begin{center}
\begin{tabular}{@{}lll@{}}
\toprule
Parameter & Value & Justification \\
\midrule
$\eta$ & $0$ & All members hash to $0^n$ by construction \\
$\varepsilon$ & $2^{-n}$ & Random $c$ satisfies $h(c\|s)=0^n$ w.p.\
    $2^{-n}$ \\
$\mu$ & $\leq 1$ & Binary output \\
$\delta$ & N/A & $K(x)=1$; no representation uniformity \\
Space & $n = -\log_2\varepsilon$ bits/element & \\
\bottomrule
\end{tabular}
\end{center}

\paragraph{Property status.}
Totality: \textbf{yes} ($h$ is defined on all inputs).
Representation uniformity: \textbf{no} ($K=1$; cipher values are just
domain elements).  Correctness: \textbf{yes} ($\eta = 0$).  Composability:
\textbf{yes} (intersection of HashSets: $x \in A \cap B$ iff
$h(x\|s_A) = 0^n$ and $h(x\|s_B) = 0^n$).

\begin{remark}
The HashSet demonstrates that the four properties are independent: a
construction can satisfy some without others.
\end{remark}

% FIXED: CM-4: clarify that "bits per element" is amortized information content
\begin{remark}[Bits per element]
The HashSet stores a single seed $s$, not per-element data.  The quantity
$\log_2(|U|/m)$ bits per element is the \emph{amortized information
content}: the total information required to specify which $m$ elements
belong to the set, divided by $m$.  This is standard in the perfect hashing
literature.
\end{remark}

\paragraph{Construction complexity.}
Each candidate seed succeeds with probability $2^{-n|A|}$ (all $|A|$
elements must hash to $0^n$), so the expected number of trials is
$2^{n|A|}$.  The space is $n|A|$ bits total, or $n = -\log_2 \varepsilon$
bits per element, matching the information-theoretic lower bound for
membership (\S\ref{subsec:lower-bound}).  The exponential construction time
motivates the two-level construction of \S\ref{sec:singular-hash}.


\subsection{Entropy Map}
\label{subsec:entropy-map}

The entropy map generalizes the HashSet to arbitrary functions
$f : X \to Y$, achieving Shannon-entropy-optimal value encoding.

\paragraph{Latent function.}
Arbitrary $f : X \to Y$ where $Y$ is finite.

\paragraph{Construction.}
Assign a prefix-free code $C_y \subseteq \{0,1\}^*$ for each $y \in Y$,
where the fraction of hash space covered by $C_y$ is proportional to
$\Pr_D[f(X) = y]$.  Find a seed $s$ (via two-level hashing) such that
\begin{equation}
h(x \| s) \in C_{f(x)} \quad \text{for all } x \in X.
\end{equation}
Evaluation: compute $h(c \| s)$ and decode via the prefix-free code.

\paragraph{Two-level hash (practical algorithm).}
\begin{enumerate}
\item A primary hash maps $x$ to a bucket index $j \in \{0, \ldots, b-1\}$.
\item For each bucket, find a local seed $s_j$ such that $h(x \| s_j)$
    decodes to $f(x)$ for all $x$ in bucket $j$.
\item Store $s_j$ in a seed table $H[j]$.
\item Query: compute bucket $j$ (1~hash), evaluate $h(x \| H[j])$ and
    decode (1~hash).  Total: $O(1)$ hash operations.
\end{enumerate}

\paragraph{Construction time.}
Let $\ell$ be the number of items in a bucket.  Each candidate seed succeeds
independently with probability $\varepsilon^\ell$ (all $\ell$ items must
land in valid codewords).  The expected number of trials per bucket is
$\varepsilon^{-\ell}$, geometric.  Total expected construction time:
$O(b \cdot \varepsilon^{-\ell})$.

\paragraph{Mapping to cipher map abstraction.}
\begin{itemize}
\item $\fhat(c) = h(c \| H[\mathrm{bucket}(c)])$, total function.
\item $\enc(x, 0) = x$ (simplest version; $K(x)=1$).
\item $\dec(r) = y$ if $r \in C_y$ for some $y$, else $\bot$.
\end{itemize}

\paragraph{Parameter instantiation.}
\begin{center}
\begin{tabular}{@{}lll@{}}
\toprule
Parameter & Value & Justification \\
\midrule
$\eta$ & Construction-dependent & Fraction of bucket failures; near-zero
    with retry \\
$\varepsilon$ & Code-dependent & Pr[random bits $\in \bigcup_y C_y$] \\
$\mu$ & $H(Y)$ & Shannon entropy of output distribution \\
$\delta$ & From code assignment & If $|C_y| \propto \Pr[f(X)=y]$,
    approaches uniform \\
Space & $-\log_2\varepsilon + H(Y)$ bits/element & \\
\bottomrule
\end{tabular}
\end{center}

\paragraph{Property status.}
Totality: \textbf{yes}.  Representation uniformity: \textbf{achievable}
(by assigning multiple codes proportional to output frequency).
Correctness: \textbf{yes} ($\eta \approx 0$, tunable).  Composability:
\textbf{yes} (output is an $n$-bit string; serves as input to another
cipher map).

\begin{proposition}[Entropy map space]
\label{prop:entropy-space}
The space per element of an entropy map for $f : X \to Y$ with output
distribution $(p_y)_{y \in Y}$ equals
$-\log_2 \varepsilon + H(Y)$
bits, where $H(Y) = -\sum_y p_y \log_2 p_y$.
\end{proposition}

\begin{proof}
The term $-\log_2 \varepsilon$ is the per-element cost of distinguishing
signal from noise: each stored element requires that its hash prefix match a
valid codeword, which occurs with probability $\varepsilon$ for random
inputs.  The term $H(Y)$ is the expected code length under
Shannon-optimal prefix-free coding, which equals the entropy of the output
distribution by Shannon's source coding theorem.  These two components are
additive because the noise-rejection bits and the value-encoding bits occupy
disjoint portions of the hash output.
\end{proof}

\begin{remark}[Connection to Bloom filters]
The entropy map for the set-indicator function $\mathbf{1}_A$ is analogous
to a Bloom filter: two codewords (``member'' / ``non-member'') allocated
hash space in ratio $\varepsilon : (1-\varepsilon)$.  The constructions
differ---hash-based lookup vs.\ bit array with $k$ hash functions---but the
space bounds are related.
\end{remark}


\subsection{Trapdoor Boolean Algebra}
\label{subsec:trapdoor-bool}

The trapdoor Boolean algebra is an online construction: no seed search is
needed, and set operations are performed via bitwise logic on hash values.

\paragraph{Latent function.}
Set operations over subsets of a string universe.  Let $X^*$ denote the free
semigroup over alphabet $X$.  The latent structure is the Boolean algebra
$(\mathcal{P}(X^*), \cap, \cup, \complement, \emptyset, X^*)$.

\paragraph{Construction.}
Define a map $F : \mathcal{P}(X^*) \to \{0,1\}^n$ by:
\begin{equation}
F(\{x_1, \ldots, x_k\}) = h(x_1) \mathbin{|} h(x_2) \mathbin{|} \cdots
\mathbin{|} h(x_k),
\end{equation}
where $|$ denotes bitwise OR and $h : X^* \to \{0,1\}^n$ is a
cryptographic hash.

The operation correspondence is:
\begin{center}
\begin{tabular}{@{}llc@{}}
\toprule
Lattice operation & Bit operation & Exact? \\
\midrule
$A \cup B$ & $F(A) \mathbin{|} F(B)$ & Yes \\
$A \cap B$ & $F(A) \mathbin{\&} F(B)$ & Approximate \\
$A^\complement$ & ${\sim}F(A)$ & Approximate \\
$\emptyset$ & $0^n$ & Yes \\
$X^*$ & $1^n$ & Yes (limiting) \\
\bottomrule
\end{tabular}
\end{center}

\paragraph{Why intersection is approximate.}
$F(A) \mathbin{\&} F(B)$ retains all bit positions set in both $F(A)$ and
$F(B)$.  But $F(A)$ has bits set by \emph{all} elements of $A$ (not just
$A \cap B$), and likewise for $F(B)$.  When elements of $A \setminus B$
happen to set the same bit positions as elements of $B \setminus A$
(cross-element hash collisions), those bits survive the AND and appear as
spurious bits.  Formally, $F(A \cap B) \subseteq F(A) \mathbin{\&} F(B)$
(as bit vectors), with equality only when there are no cross-element
collisions.

\paragraph{Why NOT degrades with set size.}
${\sim}F(A)$ flips all bits, giving the positions not set by any element of
$A$.  But $F(A^\complement)$ is the OR of hashes of all elements not in $A$:
for large complements, this converges to $1^n$.  Since $F(A)$ also
approaches $1^n$ for large $|A|$, the NOT approximation
${\sim}F(A) \approx 0^n$ is poor when $|A|$ is large.  NOT is most accurate
when $|A|$ is small relative to $n$.

\paragraph{Relational predicates.}
\begin{itemize}
\item Membership: $x \in_B W$ iff $h(x) \mathbin{\&} F(W) = h(x)$.
\item Subset: $A \subseteq_B B$ iff $F(A) \mathbin{\&} F(B) = F(A)$.
\end{itemize}

\begin{proposition}[Membership false positive rate]
\label{prop:membership-fp}
For random $x \notin W$ where $|W| = k$, under the random oracle model:
\begin{equation}
\varepsilon_{\in}(k, n) = (1 - 2^{-(k+1)})^n.
\end{equation}
\end{proposition}

\begin{proof}
Under the random oracle model, hash outputs are independent uniform bits.
Each bit position $j$ has $\Pr[h(x)_j = 1] = 1/2$ and
$\Pr[F(W)_j = 0] = (1/2)^k = 2^{-k}$ (probability none of the $k$
elements set bit $j$).  Bit $j$ causes a rejection (correctly identifying
$x \notin W$) when $h(x)_j = 1$ and $F(W)_j = 0$, which happens with
probability $2^{-(k+1)}$.  All $n$ bits must independently pass for a false
positive:
$\varepsilon_{\in} = (1 - 2^{-(k+1)})^n$.
\end{proof}

\paragraph{Space complexity.}
Single-level: $n = \log \varepsilon / \log(1 - 2^{-(k+1)}) = O(2^k)$ bits.
This is exponential in set size.  Two-level (practical): hash to $2^w$ bins
of $(q-w)$ bits each.  Total space $2^w(q-w)$ bits.  False positive rate:
$\varepsilon(k, w, q) = (1 - 2^{-(k/2^w + 1)})^{q-w}$.

\paragraph{Mapping to cipher map abstraction.}
For the membership predicate on a set $W$:
\begin{itemize}
\item $\fhat(c) = c \mathbin{\&} F(W)$, total function on all $n$-bit strings.
\item $\enc(x, 0) = h(x)$.
\item $\dec(r) = 1$ if $r = h(x)$ (where $x$ is the queried element),
    else $0$.
\end{itemize}

\paragraph{Parameter instantiation} (membership predicate, set size $k$).
\begin{center}
\begin{tabular}{@{}lll@{}}
\toprule
Parameter & Value & Justification \\
\midrule
$\eta$ & Operation-dependent & Union exact; intersection approximate \\
$\varepsilon$ & $(1-2^{-(k+1)})^n$ & Membership FP rate \\
$\mu$ & $n$ bits (fixed width) & Entire set as single $n$-bit vector \\
Space & $n$ bits total (not per element) & Constant-size; FP rate degrades
    with $k$ \\
\bottomrule
\end{tabular}
\end{center}

\paragraph{Property status.}
Totality: \textbf{yes} (bitwise operations defined on all $n$-bit strings).
Representation uniformity: \textbf{marginal only} (see
Remark~\ref{rem:marginal-only} and \S\ref{subsec:granularity}).
Correctness: \textbf{partial} (union exact; intersection/NOT approximate).
Composability: \textbf{yes} (Boolean operations compose; error propagates
via interval arithmetic).

\subsection{Construction Comparison}

\begin{table}[ht]
\centering
\caption{Comparison of the three constructions.}
\label{tab:comparison}
\begin{tabular}{@{}llll@{}}
\toprule
& HashSet & Entropy Map & Trapdoor Boolean Alg. \\
\midrule
Construction & Batch & Batch & Online \\
Latent function & $\mathbf{1}_A$ & Arbitrary $f:X\to Y$ &
    Set ops on $\mathcal{P}(X^*)$ \\
Totality & Yes & Yes & Yes \\
Rep.\ Uniformity & No ($K\!=\!1$) & Achievable & Marginal only \\
Correctness $\eta$ & $0$ & $\approx 0$ (tunable) &
    Union exact; $\cap$/$\neg$ approx. \\
$\varepsilon$ & $2^{-n}$ & Code-dependent &
    $(1-2^{-(k+1)})^n$ \\
$\mu$ & $\leq 1$ bit & $H(Y)$ & $n$ bits \\
Space/element & $-\log_2\varepsilon$ &
    $-\log_2\varepsilon + H(Y)$ & $n/k$ (amortized) \\
\bottomrule
\end{tabular}
\end{table}


%% ============================================================
\section{The Trusted/Untrusted Machine Model}
\label{sec:trust-model}

This section formalizes the trust model implicit in the constructions above.

\begin{definition}[Trusted machine $T$]
\label{def:trusted}
The trusted machine $T$ holds: the seed $s$ (trapdoor), the encoding
function $\enc$, the decoding function $\dec$, and knowledge of which
queries are real vs.\ filler.

$T$ can:
\begin{enumerate}
\item Encode plaintext values: $x \mapsto \enc(x, k)$ for chosen $k$.
\item Decode cipher results: $r \mapsto \dec(r)$.
\item Inject noise: generate random $c \gets \{0,1\}^n$ as filler queries.
\item Verify results: check $\dec(\fhat(\enc(x,k))) = f(x)$.
\item Reseed: construct a new cipher map with fresh $s'$ when leakage accumulates.
\end{enumerate}
\end{definition}

\begin{definition}[Untrusted machine $U$]
\label{def:untrusted}
The untrusted machine $U$ holds: the cipher map $\fhat$ (a total function on
bit strings) and a collection of cipher values (encodings + filler,
indistinguishable to $U$).

$U$ can:
\begin{enumerate}
\item Evaluate cipher maps: $c \mapsto \fhat(c)$ for any
    $c \in \{0,1\}^n$.
\item Return results to $T$.
\end{enumerate}

$U$ cannot:
\begin{enumerate}
\item Decode cipher values (does not hold $s$ or $\dec$).
\item Distinguish real encodings from filler (by Property~1).
\item Determine the domain $X$ or function $f$ (by Property~2, cipher
    values are $\delta$-close to uniform).
\item Enumerate which inputs are ``in-domain'' (by totality, every input
    produces output).
\end{enumerate}
\end{definition}

\paragraph{Protocol.}
The information flow is:
\begin{enumerate}
\item $T$ encodes inputs: $\enc(x_1, k_1), \ldots, \enc(x_m, k_m)$.
\item $T$ generates filler: $c_1, \ldots, c_r \gets \{0,1\}^n$.
\item $T$ sends all cipher values (real + filler, shuffled) to $U$.
\item $U$ evaluates $\fhat$ on each cipher value and returns results.
\item $T$ decodes the results it cares about; discards filler results.
\end{enumerate}

What $U$ observes is a sequence of $n$-bit strings as input and a sequence
of $n$-bit strings as output.  Under the four properties:
\begin{itemize}
\item Totality: every input produces output; no error signals leak domain
    information.
\item Representation uniformity: cipher values are $\delta$-close to
    uniform; frequency analysis is bounded by $\delta$.
\item Correctness: irrelevant to $U$ (only $T$ decodes).
\item Composability: $U$ can chain cipher maps without decoding intermediate
    results.
\end{itemize}

\begin{table}[ht]
\centering
\caption{What the four properties guarantee against $U$.}
\label{tab:guarantees}
\begin{tabular}{@{}lll@{}}
\toprule
Adversary capability & Prevented by & Mechanism \\
\midrule
Distinguish real from filler & Totality & All queries produce output \\
Frequency analysis & Rep.\ uniformity & Values $\delta$-close to uniform \\
Learn function values & Trapdoor & Cannot invert $h$ without $s$ \\
Enumerate domain $X$ & Totality + trapdoor & Cannot test membership \\
Detect query correlations & \textbf{Not fully} & Only marginal uniformity \\
\bottomrule
\end{tabular}
\end{table}

The last row is the honest limitation: the four properties do not provide
full correlation hiding unless the encoding granularity encompasses the
correlated values (\S\ref{subsec:granularity}).


%% ============================================================
\section{The Singular Hash Map}
\label{sec:singular-hash}

The Singular Hash Map is a concrete cipher map construction achieving the
information-theoretic space lower bound.  It combines the entropy map's
prefix-free coding with a global seed search.

\subsection{Information-Theoretic Lower Bound}
\label{subsec:lower-bound}

\begin{theorem}[Lower bound]
\label{thm:lower-bound}
Any data structure representing an approximate map $\fhat$ for
$f : X \to Y$ with noise decode probability $\varepsilon$ and correctness
$\eta = 0$ over $n$ stored elements requires at least
\begin{equation}
n\bigl(-\log_2 \varepsilon + H(Y)\bigr) \text{ bits},
\end{equation}
or equivalently $-\log_2 \varepsilon + H(Y)$ bits per element.
\end{theorem}

\begin{proof}
The bound decomposes into two independent information requirements.

% FIXED: CM-1: replaced incorrect Bloom citation with self-contained counting argument
\emph{Membership component.}
Consider the membership predicate $\dec(\fhat(c)) \neq \bot$.  For a
universe $U$ with $|U| \gg n$, the structure must distinguish $n$ stored
elements from $|U|-n$ non-elements, with at most an $\varepsilon$ fraction
of non-elements decoding to valid output.  There are $\binom{|U|}{n}$
possible sets of size $n$ from a universe $U$.  Any representation must
distinguish all such sets, requiring at least $\log_2 \binom{|U|}{n}$
bits.  For $|U| \gg n$, Stirling's approximation gives
$\log_2 \binom{|U|}{n} \approx n \log_2(|U|/n)$ bits, or
$\log_2(|U|/n)$ bits per element.  Setting $\varepsilon = n/|U|$ (the
fraction of the universe that decodes validly), this is
$n \log_2(1/\varepsilon)$ bits.

\emph{Value component.}
Independently of membership, the structure must encode the function values
$(f(x_1), \ldots, f(x_n))$.  By Shannon's source coding theorem, encoding
$n$ independent draws from the distribution $(p_y)$ requires at least
$n H(Y)$ bits.

% FIXED: Added justification for additivity (entropy chain rule / independence)
Since membership and value encoding are jointly necessary, the total
requirement is $n(-\log_2 \varepsilon + H(Y))$ bits.  The additivity holds
because the membership bits and value-encoding bits convey independent
information: by the entropy chain rule,
$H(\text{membership}, \text{value}) = H(\text{membership}) + H(\text{value} \mid \text{membership})$,
and conditioning on membership (i.e., knowing which elements are stored) does
not reduce the entropy of the value assignment, since the values
$(f(x_1), \ldots, f(x_n))$ are an independent degree of freedom.
\end{proof}

\subsection{Construction}

\begin{algorithm}
\caption{Singular Hash Map Construction}
\label{alg:singular-hash}
\begin{algorithmic}[1]
\Function{BuildCipherMap}{$M$, $\enc$, $\dec$, $h$}
    \State $\ell \gets 0$
    \While{true}
        \State $\text{success} \gets \text{true}$
        \For{$(x, y) \in M$}
            \State $\text{hash} \gets h(\ell) \oplus h(x)$
                \Comment{XOR combines seed with key}
            \If{$\dec(\text{hash}) \neq y$}
                \State $\text{success} \gets \text{false}$; \textbf{break}
            \EndIf
        \EndFor
        \If{success}
            \State \Return $\text{CipherMap}(h, \dec, \ell)$
        \EndIf
        \State $\ell \gets \ell + 1$
    \EndWhile
\EndFunction
\end{algorithmic}
\end{algorithm}

The construction searches for a seed $\ell$ such that for every
$(x, y) \in M$, the hash $h(\ell) \oplus h(x)$ decodes to $y$ under the
prefix-free code.  Query evaluation computes
$\fhat(c) = h(\ell) \oplus h(c)$ and decodes the result.

% FIXED: CM-2: clarify that Algorithm 1 implements eta=0; describe eta>0 relaxation
\begin{remark}[Construction with $\eta > 0$]
Algorithm~\ref{alg:singular-hash} implements the $\eta = 0$ special case.
For $\eta > 0$, the search accepts any seed for which at most
$\lfloor \eta \cdot m \rfloor$ elements fail correctness.  This relaxation
reduces expected construction time since fewer constraints must be
simultaneously satisfied.
\end{remark}

\subsection{Space Optimality}

\begin{theorem}[Space complexity]
\label{thm:space-optimal}
The Singular Hash Map achieves space complexity
\begin{equation}
(1-\eta)\bigl(-\log_2 \varepsilon + \mu\bigr) \text{ bits per element},
\end{equation}
where $\mu = H(Y)$ is the mean encoding length.  When $\eta = 0$, this
matches the lower bound of Theorem~\ref{thm:lower-bound}.
\end{theorem}

\begin{proof}
\emph{Step 1: False negative reduction.}
When false negatives are permitted at rate $\eta$, only $(1-\eta)n$ of the
$n$ elements must decode correctly, reducing the effective element count.

\emph{Step 2: Per-element bit cost.}
For each stored element $x_i$, the hash $h(\ell) \oplus h(x_i)$ must decode
to $f(x_i)$ under the prefix-free code.  Under the random oracle model, the
probability of correct decode for element $x_i$ with target value $y_i$ is
$2^{-|c_{y_i}|}$, where $|c_{y_i}|$ is the code length.  By
Shannon-optimal coding, $|c_{y_i}| \approx -\log_2 p_{y_i}$, so the
expected code length per element is
$\sum_y p_y |c_y| = \mu = H(Y)$.

\emph{Step 3: False positive control.}
For non-stored elements $x \notin M$, the hash $h(\ell) \oplus h(x)$ is
uniformly distributed under the random oracle model.  The probability of
decoding to \emph{any} valid output is $\varepsilon$ (the total probability
mass of the prefix-free code).  Achieving target rate $\varepsilon$ requires
$-\log_2 \varepsilon$ bits of hash beyond the value encoding.

\emph{Step 4: Combining.}
Each stored element costs $-\log_2 \varepsilon + \mu$ bits, and
$(1-\eta)n$ elements are stored, giving total space
$(1-\eta)n(-\log_2 \varepsilon + \mu)$ bits.
\end{proof}

\begin{corollary}[Asymptotic optimality]
When $\eta = 0$, the Singular Hash Map achieves $-\log_2 \varepsilon + H(Y)$
bits per element, matching the information-theoretic lower bound.
\end{corollary}

\subsection{Construction Time}

\begin{proposition}[Expected construction trials]
\label{prop:construction-time}
For $m = (1-\eta)n$ stored elements, the probability of successful
construction for a given seed $\ell$ is:
\begin{equation}
p = \prod_{i=1}^{m} (1 - e_i),
\end{equation}
where $e_i = 1 - 2^{-|c_{y_i}|}$ is the per-element failure probability.
The expected number of trials is $1/p$.
\end{proposition}

\begin{proof}
Under the random oracle model, the hash values $h(\ell) \oplus h(x_i)$ for
distinct $x_i$ are independent uniformly distributed bit strings.  For
element $x_i$ with target value $y_i$, correct decoding requires the hash
prefix to match code $c_{y_i}$, which occurs with probability
$2^{-|c_{y_i}|}$.  Independence across elements gives the product formula.
The number of trials is geometric with success probability $p$.
\end{proof}


%% ============================================================
\section{Composition}
\label{sec:composition}

\subsection{Warm-Up: The AND Gate}
\label{subsec:and-gate}

Before proving the general composition theorem, we analyze a concrete case:
the AND gate operating on two independent Bernoulli Boolean values
$B_1, B_2$ with correctness probabilities $p_1 = 1-\eta_1$ and
$p_2 = 1-\eta_2$.

\begin{proposition}[AND gate correctness]
\label{prop:and-gate}
The correctness probability of $\mathrm{AND}(B_1, B_2)$ as an approximation
of $\mathrm{AND}(x_1, x_2)$ depends on the latent inputs:
\begin{center}
\begin{tabular}{@{}cccc@{}}
\toprule
$x_1$ & $x_2$ & $\mathrm{AND}(x_1,x_2)$ & $\Pr[\text{correct}]$ \\
\midrule
$1$ & $1$ & $1$ & $p_1 p_2$ \\
$1$ & $0$ & $0$ & $1 - p_1(1-p_2)$ \\
$0$ & $1$ & $0$ & $1 - p_2(1-p_1)$ \\
$0$ & $0$ & $0$ & $p_1 + p_2 - p_1 p_2$ \\
\bottomrule
\end{tabular}
\end{center}
\end{proposition}

\begin{proof}
We verify the $(1,0)$ case; the others are analogous.  When $x_1=1$,
$x_2=0$, the correct output is $0$.  The output is wrong ($=1$) only when
$B_1=1$ (correct, probability $p_1$) \emph{and} $B_2=1$ (incorrect,
probability $1-p_2$).  So
$\Pr[\text{wrong}] = p_1(1-p_2)$ and
$\Pr[\text{correct}] = 1 - p_1(1-p_2) = 1 - p_1 + p_1 p_2$.
\end{proof}

The output is a \emph{4th-order Bernoulli Boolean}: four distinct
correctness probabilities despite the Boolean type having only two values.
Rather than tracking all four, we can bound them by an interval
$[\min_{\text{cases}} p_{\text{correct}},\; \max_{\text{cases}} p_{\text{correct}}]$.
For AND with $p_1, p_2 \in (0.5, 1)$: minimum $p_1 p_2$ (case $(1,1)$);
maximum $p_1 + p_2 - p_1 p_2$ (case $(0,0)$).

\subsection{General Composition Theorem}

\begin{theorem}[Composition correctness]
\label{thm:composition}
Let $\fhat$ be a cipher map for $f : X \to Y$ with correctness $\eta_f$,
and $\ghat$ a cipher map for $g : Y \to Z$ with correctness $\eta_g$.  Then
$\ghat \circ \fhat$ is a cipher map for $g \circ f$ with correctness:
\begin{equation}
\eta_{g \circ f} = 1 - (1 - \eta_f)(1 - \eta_g) = \eta_f + \eta_g -
\eta_f \eta_g.
\end{equation}
\end{theorem}

\begin{proof}
The composition is correct when \emph{both} maps produce correct results.
We establish two bounds.

\emph{Upper bound (union bound).}
\begin{equation}
\Pr[\text{incorrect}] \leq \Pr[\fhat \text{ incorrect}] +
    \Pr[\ghat \text{ incorrect}] = \eta_f + \eta_g.
\end{equation}

\emph{Tight bound (independence).}
Assume: (i)~$\fhat$ and $\ghat$ use independent seeds, and
(ii)~$\ghat$'s correctness on $\fhat(c)$ is independent of whether $\fhat$
was correct (re-randomization: the encoding fed to $\ghat$ is a fresh
random valid encoding of $f(x)$).  Then
\begin{equation}
\Pr[\text{both correct}] = (1-\eta_f)(1-\eta_g),
\end{equation}
giving
$\eta_{g \circ f} = 1 - (1-\eta_f)(1-\eta_g) = \eta_f + \eta_g -
\eta_f \eta_g$.
\end{proof}

\begin{remark}[Independence assumption]
Condition~(ii) (re-randomization) deserves care.  In practice, $\fhat$
produces a \emph{specific} encoding of $f(x)$, not a uniformly random one.
If $\ghat$'s correctness depends on \emph{which} encoding it receives, the
errors may be correlated.  The formula is therefore best understood as an
upper bound under the weaker assumption that errors are positively
correlated.  For specific circuits, the case-by-case analysis from
\S\ref{subsec:and-gate} (interval arithmetic) gives tighter bounds.
\end{remark}

% FIXED: CM-3: address garbage propagation when f-hat fails
\begin{remark}[Garbage propagation]
When $\fhat$ produces an incorrect output on some input, the value fed to
$\ghat$ is effectively random.  Since $\ghat$ is a total function, it
produces \emph{some} output on this random input, which is correct with
probability at most $\varepsilon_g$ (the noise-decode probability of
$\ghat$).  Thus the bound
$\eta_{\mathrm{total}} \leq \eta_f + \eta_g - \eta_f \eta_g$ is an upper
bound; the actual error rate may be marginally lower due to accidental
correctness on garbage inputs.  For small $\eta_f, \eta_g$, this
distinction is negligible.
\end{remark}

\subsection{Composition Chains}

\begin{corollary}[Chain composition]
\label{cor:chain}
For a chain of $m$ cipher maps $\fhat_1, \ldots, \fhat_m$:
\begin{equation}
\eta_{\mathrm{total}} = 1 - \prod_{i=1}^{m}(1 - \eta_i).
\end{equation}
If all $\eta_i = \eta$:
\begin{equation}
\eta_{\mathrm{total}} = 1 - (1-\eta)^m \approx m\eta
\quad \text{for small } \eta.
\end{equation}
\end{corollary}

\begin{proof}
By induction on Theorem~\ref{thm:composition}.  The base case $m=2$ is the
theorem.  For the inductive step, the composition of $m-1$ maps has
correctness $1 - \prod_{i=1}^{m-1}(1-\eta_i)$; composing with $\fhat_m$
gives
$1 - \prod_{i=1}^{m-1}(1-\eta_i) \cdot (1-\eta_m) = 1 - \prod_{i=1}^m
(1-\eta_i)$.
The approximation follows from $\ln(1-\eta) \approx -\eta$ for small
$\eta$, giving
$(1-\eta)^m \approx e^{-m\eta} \approx 1 - m\eta$.
\end{proof}

\begin{example}
A circuit of 100 cipher maps each with $\eta = 10^{-6}$:
$\eta_{\mathrm{total}} = 1 - (1-10^{-6})^{100} \approx 10^{-4}$.
\end{example}

\subsection{Error Accumulation by Gate Type}

For specific circuits, the AND gate analysis
(Proposition~\ref{prop:and-gate}) can be extended to OR, NOT, and arbitrary
Boolean gates.  Each gate type induces a different case table.  Rather than
tracking exact correctness per case, interval arithmetic propagates bounds:
at each gate, compute $[\eta_{\min}, \eta_{\max}]$ from the input intervals.
This gives input-dependent bounds tighter than the worst-case
Theorem~\ref{thm:composition}.


%% ============================================================
\section{Representation Uniformity and Encoding Granularity}
\label{sec:uniformity}

\subsection{Total Variation Distance}

\begin{definition}[Total variation distance]
\label{def:tv-distance}
Let $P$ and $Q$ be distributions on a finite set $\Omega$.  The total
variation distance is:
\begin{equation}
\TV(P, Q) = \frac{1}{2} \sum_{\omega \in \Omega} |P(\omega) - Q(\omega)|
= \max_{A \subseteq \Omega} |P(A) - Q(A)|.
\end{equation}
\end{definition}

TV distance has the operational interpretation: no test can distinguish $P$
from $Q$ with advantage greater than $\TV(P, Q)$.  This directly bounds
what an adversary observing cipher values can learn.

\begin{remark}
Max-divergence ($D_\infty$) would give a stronger guarantee (bounding the
\emph{ratio} of probabilities rather than the difference) but is harder to
achieve.  We use TV distance because the constructions are designed around
frequency equalization, which naturally yields TV bounds.
\end{remark}

\subsection{The Encoding Granularity Principle}
\label{subsec:granularity}

Representation uniformity is defined \emph{per cipher map}.  What
correlations it hides depends on what is encoded as a unit.

\begin{definition}[Encoding granularity]
The \emph{encoding granularity} of a cipher map $\fhat$ for $f : X \to Y$
is the type $X$.  Representation uniformity applies to the distribution over
encodings of individual elements of $X$.
\end{definition}

Consider encrypting a pair of Boolean values $(a, b)$ where $a$ and $b$ are
correlated.

\begin{description}
\item[Component-wise encoding ($p=1$).]
Encode $a$ and $b$ as separate cipher values.  Each component's marginal
distribution is $\delta$-close to uniform, but joint correlations leak.

\item[Pair-wise encoding ($p=2$).]
Encode $(a, b)$ as a single cipher value.  The distribution over pair
encodings is $\delta$-close to uniform; correlations between $a$ and $b$ are
hidden.  Cost: domain is $\{0,1\}^2$ (4 elements) rather than $\{0,1\}$
(2 elements); space per encoding grows.

\item[Whole-state encoding ($p=k$).]
Encode the entire program state as one cipher value.  All correlations are
hidden.  Cost: domain is the full state space, typically prohibitive.
\end{description}

\begin{definition}[Entanglement parameter]
For a system encoding $k$ correlated Boolean values, the \emph{entanglement
parameter} $p$ is the number of values encoded as a single unit.  The system
has type $\mathrm{cipher}(\{0,1\}^p)^{k/p}$, where each block of $p$ bits
is encoded jointly.
\end{definition}

\begin{proposition}[Granularity and privacy]
\label{prop:granularity}
Let $\fhat$ be a cipher map for $f : X_1 \times X_2 \to Y$ with
representation uniformity $\delta$.  Then the marginal distribution of
$\enc((x_1, x_2), k)$ over random $(x_1, x_2) \sim D$ and random $k$ is
$\delta$-close to uniform, and the joint distribution of $(x_1, x_2)$ is
hidden up to $\delta$.

Conversely, let $\fhat_1, \fhat_2$ be independent cipher maps for
$f_1 : X_1 \to Y_1$ and $f_2 : X_2 \to Y_2$, each with representation
uniformity $\delta$.  Then each marginal cipher value is $\delta$-close to
uniform, but correlations between $x_1$ and $x_2$ are \textbf{not hidden},
even when $\delta = 0$.
\end{proposition}

\begin{proof}
The first part follows directly from Definition~\ref{def:rep-uniformity}
applied to the joint cipher map.  For the second part: even with $\delta=0$
(perfect marginal uniformity), $\enc_i$ is a deterministic function of
$(x_i, k_i)$.  An adversary observing $N$ draws of
$(\enc_1(x_1, k_1), \enc_2(x_2, k_2))$ can estimate the joint distribution
and recover the correlation structure of $(x_1, x_2)$, because the
correlation structure under $D$ is preserved in the joint distribution of
cipher value pairs.
\end{proof}

\paragraph{The trade-off.}
Coarser encoding granularity (larger $p$) yields better privacy but
requires more space and construction time.  Composition exists precisely
because encoding the entire program as a single cipher value is impractical:
we compose smaller cipher maps and accept that intermediate correlations may
leak.

\paragraph{Honest limitations.}
\begin{enumerate}
\item Marginal uniformity says nothing about joint distributions.
\item Encoding granularity is a design choice, not a property of the
    formalism.
\item Equality pattern leakage is fundamental: any deterministic encoding
    leaks the equality pattern of its inputs.  Multiple representations
    ($K(x) > 1$) reduce but do not eliminate this.
\end{enumerate}


%% ============================================================
\section{Discussion}
\label{sec:discussion}

\subsection{Relationship to the Bernoulli Model}

The correctness parameter $\eta$ of a cipher map corresponds directly to the
error rate in a Bernoulli approximation.  A cipher map with correctness
$\eta$ induces a Bernoulli approximation of the latent function $f$: for
each $x \in X$,
$\Pr[\dec(\fhat(\enc(x,k))) \neq f(x)] = \eta(x)$,
where the failing elements are deterministic for a given seed.

The composition formula
$\eta_{g \circ f} = 1 - (1-\eta_f)(1-\eta_g)$
is the standard error accumulation for independent Bernoulli events,
identical to the formula derived in the noisy gates framework for AND gates
over Bernoulli Booleans.  This is not coincidental: a cipher map \emph{is} a
Bernoulli approximation, viewed through the lens of hash-based total
functions.

\subsection{Algebraic Structure}

The cipher map abstraction relates to algebraic concepts from the theory of
monoid homomorphisms.  The decoding function $\dec$ is (by construction) a
left inverse of $\enc$, and for cipher maps that preserve algebraic
structure (e.g., the trapdoor Boolean algebra), $\dec$ is a homomorphism
from the cipher monoid to the latent monoid---but only up to the
equivalence $c_1 \sim c_2 \iff \dec(c_1) = \dec(c_2)$.  The quotient
$\{0,1\}^n / {\sim}$ is isomorphic to the latent structure.

We note this connection but do not develop it further.  The category-theoretic
framing (cipher functors lifting monoids to cipher monoids) is a natural
extension but requires careful treatment of the approximate setting: the
functor laws hold exactly on the quotient but only approximately on
representatives.

\subsection{Open Questions}

\begin{enumerate}
\item \textbf{Max-divergence vs.\ TV distance for $\delta$.}
    TV distance bounds distinguishing advantage; max-divergence
    ($D_\infty$) would give per-element bounds.  Which is the right metric
    depends on the threat model.

\item \textbf{Tight bounds for intersection and NOT.}
    The intersection false positive rate from cross-element collisions in
    the trapdoor Boolean algebra needs a closed-form expression as a
    function of $|A|$, $|B|$, $|A \cap B|$, and $n$.  For NOT, deriving
    $\eta_{\mathrm{NOT}}(|A|, n)$ would complete the parameter instantiation.

\item \textbf{Adaptive encoding granularity.}
    Can the entanglement parameter $p$ be made adaptive---encoding pairs
    only for correlated values while encoding independent values
    separately?

\item \textbf{Noise-to-signal probability under composition.}
    Through a chain of $m$ cipher maps, what is the probability that a
    filler query produces a decodable output at the end?  Naively
    $1 - (1-\varepsilon)^m$, but this assumes independence of the
    $\varepsilon$ events across maps.

\item \textbf{Layered construction laws.}
    Do the three construction layers (undef, noise, cipher) satisfy formal
    algebraic laws in the approximate setting?  If so, this would enable
    equational reasoning about cipher map constructions.
\end{enumerate}

\subsection{What This Framework Is Not}

To set clear expectations:
\begin{itemize}
\item \emph{Not ORAM.}  ORAM hides access patterns by reshuffling storage
    with polylogarithmic overhead.  Cipher maps use static total functions.
\item \emph{Not differential privacy.}  Differential privacy bounds what
    can be learned about a \emph{dataset} from a \emph{query answer}.
    Cipher maps bound what an untrusted evaluator learns about a
    \emph{function} from \emph{opaque evaluations}.
\item \emph{Not simulation-based security.}  We do not claim that the
    untrusted machine's view can be simulated.  The guarantees are
    parameterized ($\delta$, $\varepsilon$) rather than negligible in a
    security parameter.
\end{itemize}


%% ============================================================
\section*{Acknowledgments}

The constructions in this paper draw on the author's earlier work on the
Bernoulli data type library and the trapdoor Boolean algebra.

\bibliographystyle{plainnat}
\bibliography{references}

\end{document}
